% BEGIN REV09_RECTAS27_PROLONGACION
\subsubsection{Prolongement réversible en coordonnées de droites}
\label{corpus9:xii:rectas-prolongacion}

La proposition~\ref{prop:cobertura-729-seccion} a construit les
bijections \(\beta_{{\rm blk},N}\) au moyen de segments nonadiques du
registre ordonné, avec leur compatibilité avec la suppression des
blocs. L’homéomorphisme
\eqref{eq:homeomorfismo-calibrado-ch} identifie la limite de ces
histoires engendrées à la suite de leurs mots de six trits.
Nous composons maintenant ce lecteur avec le système de coordonnées fini
\eqref{corpus9:xii:rectas-codificacion}, en conservant sa loi centrale.

\begin{theorem}[Transport cylindrique des histoires vers des paires de droites]
\label{corpus9:xii:rectas-homeomorfismo}
Pour chaque \(N\geq0\), l’application
\begin{equation}
 \mathcal E_N:=\mathcal E^{\times N}\circ\beta_{{\rm blk},N}:
 Y_N\xrightarrow{\ \cong\ }
 (\mathscr L_{27}\times\mathscr L_{27})^N
 \label{corpus9:xii:rectas-nivel}
\end{equation}
est une bijection. Ces bijections satisfont
\begin{equation}
 \operatorname{pref}_N\circ\mathcal E_{N+1}
 =\mathcal E_N\circ\rho^{\rm blk}_{N+1,N}
 \label{corpus9:xii:rectas-truncamiento}
\end{equation}
et déterminent l’homéomorphisme

\begin{equation}
 \mathcal E_\infty:=\mathcal E^{\mathbb N}\circ\beta_{\rm blk}:
 \Omega_\Gamma^{\rm cal}\xrightarrow{\ \cong\ }
 (\mathscr L_{27}\times\mathscr L_{27})^{\mathbb N}.
 \label{corpus9:xii:rectas-limite}
\end{equation}
Son inverse est
\(\beta_{\rm blk}^{-1}\circ(\mathcal E^{-1})^{\mathbb N}\).

\end{theorem}
\begin{proof}
La proposition~\ref{corpus9:xii:rectas-biyeccion} fournit
l’inverse effectif de chaque facteur \(\mathcal E\), au moyen du tableau
des droites et de l’entrelacement des coordonnées. C’est pourquoi
\(\mathcal E_N^{-1}=\beta_{{\rm blk},N}^{-1}
\circ(\mathcal E^{-1})^{\times N}\). Pour \(N=0\), les deux applications
agissent sur le mot vide. L’application coordonnée par coordonnée
commute avec la suppression du dernier facteur. En utilisant
\eqref{eq:cobertura-729-ch}, on obtient
\[
 \begin{aligned}
 \operatorname{pref}_N\mathcal E^{\times(N+1)}
           \beta_{{\rm blk},N+1}
 &=\mathcal E^{\times N}\operatorname{pref}_N
           \beta_{{\rm blk},N+1}\\
 &=\mathcal E^{\times N}\beta_{{\rm blk},N}
           \rho^{\rm blk}_{N+1,N}.
 \end{aligned}
\]
C’est \eqref{corpus9:xii:rectas-truncamiento}. Une suite de paires
de droites détermine, par l’inverse de chaque facteur, une suite unique
de mots. L’homéomorphisme
\eqref{eq:homeomorfismo-calibrado-ch} reconstruit à partir de celle-ci une
histoire unique dans \(\Omega_\Gamma^{\rm cal}\). Les compositions
récupèrent respectivement l’histoire et la suite initiales.
Dans les topologies cylindriques, un ouvert de base fixe un nombre fini
de coordonnées. L’application et son inverse transforment ces conditions
en conditions sur le même nombre de coordonnées ; toutes deux sont
continues. Cela démontre l’affirmation topologique et la compatibilité
de la limite avec toutes les troncatures finies.

\end{proof}

Pour conserver l’information de l’état complet, nous notons
\(\zeta\) ses registres restants : repère, orientation, retenue,
chemin, mémoire et données de frontière. Sur le domaine admissible déjà
construit, la transformation étendue est
\[
 \widehat{\mathcal E}(w,\zeta)=(\mathcal E(w),\zeta).
\]
Son domaine image conserve les mêmes conditions de compatibilité,
exprimées au moyen de l’inverse explicite de la première coordonnée. Si
\(F:\mathcal X\to\mathcal Y\) est un opérateur entre domaines d’états,
son transport est
\[
 F^{\mathcal E}
 =\widehat{\mathcal E}_{\mathcal Y}\circ F
       \circ\widehat{\mathcal E}_{\mathcal X}^{-1}.
\]
L’identité \(F^{\mathcal E}\widehat{\mathcal E}_{\mathcal X}
=\widehat{\mathcal E}_{\mathcal Y}F\) s’obtient en annulant les deux
applications inverses. Pour un prolongement relationnel
\(\mathcal R\subset\mathcal X\times\mathcal Y\), la définition
correspondante est
\[
 \mathcal R^{\mathcal E}
 =\{(\widehat{\mathcal E}_{\mathcal X}x,
       \widehat{\mathcal E}_{\mathcal Y}y):(x,y)\in\mathcal R\}.
\]
La transformation inverse récupère exactement \(\mathcal R\).
Si le prolongement conserve en outre un témoin \(\omega\) dans un domaine
\(W\), sa réalisation \(\widetilde{\mathcal R}\subset
\mathcal X\times W\times\mathcal Y\) est transportée par
\[
 (x,\omega,y)\longmapsto
 (\widehat{\mathcal E}_{\mathcal X}x,\omega,
       \widehat{\mathcal E}_{\mathcal Y}y).
\]
L’inverse agit sur les deux extrémités et conserve \(\omega\). Ainsi,
des témoins différents restent différents, même s’ils partagent les
extrémités, et la multiplicité enregistrée demeure. La condition de
fonctionnalité se conserve sur les domaines où elle était déjà établie.

La couverture à profondeur arbitraire provient ici du domaine
\(\Omega_\Gamma^{\rm cal}\) et de son lecteur démontré, tandis que les
paires de droites fournissent des coordonnées réversibles de ce domaine.
Le recensement initial de \(243\) mots et la capacité globale de
\(729\) prolongements par bloc conservent leurs fonctions distinctes.
La classification diagonale, incidente et gauchie accompagne chaque facteur ;
les conditions qui sélectionnent des marches adjacentes se transportent comme
restrictions de la suite complète.

% END REV09_RECTAS27_PROLONGACION
